\documentclass[10pt]{extarticle}
\usepackage[a4paper,margin=1.5cm]{geometry}
\usepackage[T1]{fontenc}
\usepackage{amsmath,amssymb}
\usepackage{graphicx}
\usepackage[export]{adjustbox}
\usepackage{multicol}
\usepackage{enumitem}
\usepackage{booktabs,array,tabularx}
\setlength{\columnsep}{0.7cm}
\setlength{\parindent}{0pt}
\setlength{\parskip}{3pt}
\pagestyle{empty}
\makeatletter
% figure registry (figdefs.tex is generated by the extract step)
\newcommand{\qfigdef}[4]{\@namedef{qf@file@#1}{#2}\@namedef{qf@w@#1}{#3}\@namedef{qf@wide@#1}{#4}}
\newcommand{\qf@check}[1]{\@ifundefined{qf@file@#1}{\PackageError{score2}{Unknown figure id #1}{}}{}}
\newcommand{\qf@img}[4][]{\includegraphics[width=#2\dimexpr\@nameuse{qf@w@#4}pt\relax,max width=#3,max height=0.45\textheight,#1]{\@nameuse{qf@file@#4}}}
% Full-width questions: finalize wraps every question that contains a wide figure in \qfullbegin...\qfullend,
% so the WHOLE question (stem, figures, options) is set across the page in reading order, then the two columns
% resume. The question's head - its number, stem and first figure (\qfullsplit marks the end of the head) -
% is kept on one page, so a number is never left at a page foot with its figure overleaf; the rest may break
% across pages normally, so no big gaps are left. Without a \qfullsplit, the whole question is kept together
% if it is under 45% of a page. Never write these yourself.
\newif\ifqf@full
\newif\ifqf@boxed
\newsavebox{\qf@box}
\newcommand{\qf@open}{\global\qf@fulltrue\par\global\qf@boxedtrue
  \setbox\qf@box\vbox\bgroup\hsize\textwidth\linewidth\textwidth}
\newcommand{\qf@emit}[1]{\par\egroup\global\qf@boxedfalse
  \ifdim\dimexpr\ht\qf@box+\dp\qf@box\relax<#1\textheight\noindent\usebox{\qf@box}\par\else\unvbox\qf@box\fi}
\newcommand{\qfullsplit}{\ifqf@boxed\qf@emit{0.85}\fi}
\newcommand{\qf@close}{\ifqf@boxed\qf@emit{0.45}\fi\global\qf@fullfalse}
\newcommand{\qfullbegin}{\end{multicols}\qf@open}
\newcommand{\qfullend}{\qf@close\begin{multicols}{2}}
\newcommand{\qfullstart}{\qf@open}
\newcommand{\qfulldone}{\qf@close}
% \qfig[scale]{id}: block figure, centred, directly where it is written (size set by the tool).
% Fallback if a wide figure is not inside a full-width question: break out of the columns in place - no floats.
\newcommand{\qfig}[2][1]{\qf@check{#2}%
  \ifqf@full\qf@blk{#1}{#2}\else
  \ifnum\@nameuse{qf@wide@#2}=1 \expandafter\@firstoftwo\else\expandafter\@secondoftwo\fi
  {\end{multicols}\begin{center}\qf@img{#1}{\textwidth}{#2}\end{center}\begin{multicols}{2}}%
  {\qf@blk{#1}{#2}}\fi}
\newcommand{\qf@blk}[2]{\par\smallskip{\centering\qf@img{#1}{\linewidth}{#2}\par}\smallskip}
% \qfigi[scale]{id}: inline figure (inside option lists / tables / minipages); never breaks columns.
\newcommand{\qfigi}[2][1]{\qf@check{#2}\qf@img[valign=c]{#1}{\linewidth}{#2}}
\makeatother
% structure macros - lint reads these to check completeness and figure placement
\newcommand{\papersec}[1]{\par\bigskip{\centering\Large\bfseries PAPER-#1\par}\medskip}
\newcommand{\partsec}[2]{\par\medskip\noindent{\large\bfseries PART-#1: #2}\par\smallskip}
\newcommand{\sect}[1]{\par\smallskip\noindent\textbf{#1}\par}
\newcommand{\qn}[1]{\par\smallskip\noindent\textbf{#1.}~}
\newcommand{\paraq}[2]{\par\smallskip\noindent\textbf{Paragraph for Questions #1 and #2}\par}
\newenvironment{qoptions}{\begin{itemize}[leftmargin=1.8em,itemsep=1pt,topsep=2pt,label={}]}{\end{itemize}}
\newcommand{\opt}[2]{\item[(#1)] #2}
% figure registry (generated by score2_tools.py finalize)
\qfigdef{p1_phy_s1_q01_1}{Score Advanced - Test 19 (04 Apr 2024)-img/p1_phy_s1_q01_1.png}{74.9}{0}
\qfigdef{p1_phy_s1_q02_1}{Score Advanced - Test 19 (04 Apr 2024)-img/p1_phy_s1_q02_1.png}{199.4}{0}
\qfigdef{p1_phy_s2_q06_1}{Score Advanced - Test 19 (04 Apr 2024)-img/p1_phy_s2_q06_1.png}{72.3}{0}
\qfigdef{p1_phy_s2_q07_1}{Score Advanced - Test 19 (04 Apr 2024)-img/p1_phy_s2_q07_1.png}{74.2}{0}
\qfigdef{p1_phy_s2_q08_1}{Score Advanced - Test 19 (04 Apr 2024)-img/p1_phy_s2_q08_1.png}{118.2}{0}
\qfigdef{p1_phy_s2_q09_1}{Score Advanced - Test 19 (04 Apr 2024)-img/p1_phy_s2_q09_1.png}{237.6}{0}
\qfigdef{p1_phy_s2_q10_1}{Score Advanced - Test 19 (04 Apr 2024)-img/p1_phy_s2_q10_1.png}{144.9}{0}
\qfigdef{p1_phy_s2_q10_A}{Score Advanced - Test 19 (04 Apr 2024)-img/p1_phy_s2_q10_A.png}{144.6}{0}
\qfigdef{p1_phy_s2_q10_B}{Score Advanced - Test 19 (04 Apr 2024)-img/p1_phy_s2_q10_B.png}{145.9}{0}
\qfigdef{p1_phy_s3_q01_1}{Score Advanced - Test 19 (04 Apr 2024)-img/p1_phy_s3_q01_1.png}{94.2}{0}
\qfigdef{p1_phy_s3_q02_1}{Score Advanced - Test 19 (04 Apr 2024)-img/p1_phy_s3_q02_1.png}{81.5}{0}
\qfigdef{p1_phy_s3_q04_1}{Score Advanced - Test 19 (04 Apr 2024)-img/p1_phy_s3_q04_1.png}{62.5}{0}
\qfigdef{p1_phy_s3_q05_1}{Score Advanced - Test 19 (04 Apr 2024)-img/p1_phy_s3_q05_1.png}{84.6}{0}
\qfigdef{p1_chem_s1_q03_1}{Score Advanced - Test 19 (04 Apr 2024)-img/p1_chem_s1_q03_1.png}{112.3}{0}
\qfigdef{p1_chem_s1_q03_A}{Score Advanced - Test 19 (04 Apr 2024)-img/p1_chem_s1_q03_A.png}{54.6}{0}
\qfigdef{p1_chem_s1_q03_B}{Score Advanced - Test 19 (04 Apr 2024)-img/p1_chem_s1_q03_B.png}{53.6}{0}
\qfigdef{p1_chem_s1_q03_C}{Score Advanced - Test 19 (04 Apr 2024)-img/p1_chem_s1_q03_C.png}{65.4}{0}
\qfigdef{p1_chem_s1_q03_D}{Score Advanced - Test 19 (04 Apr 2024)-img/p1_chem_s1_q03_D.png}{71.8}{0}
\qfigdef{p1_chem_s1_q04_1}{Score Advanced - Test 19 (04 Apr 2024)-img/p1_chem_s1_q04_1.png}{237.1}{0}
\qfigdef{p1_chem_s1_q04_A}{Score Advanced - Test 19 (04 Apr 2024)-img/p1_chem_s1_q04_A.png}{46.7}{0}
\qfigdef{p1_chem_s1_q04_B}{Score Advanced - Test 19 (04 Apr 2024)-img/p1_chem_s1_q04_B.png}{65.9}{0}
\qfigdef{p1_chem_s1_q04_C}{Score Advanced - Test 19 (04 Apr 2024)-img/p1_chem_s1_q04_C.png}{62.2}{0}
\qfigdef{p1_chem_s1_q04_D}{Score Advanced - Test 19 (04 Apr 2024)-img/p1_chem_s1_q04_D.png}{65.1}{0}
\qfigdef{p1_chem_s2_q07_1}{Score Advanced - Test 19 (04 Apr 2024)-img/p1_chem_s2_q07_1.png}{276.3}{1}
\qfigdef{p1_chem_s2_q09_1}{Score Advanced - Test 19 (04 Apr 2024)-img/p1_chem_s2_q09_1.png}{245.2}{0}
\qfigdef{p1_chem_s2_q09_A}{Score Advanced - Test 19 (04 Apr 2024)-img/p1_chem_s2_q09_A.png}{46.1}{0}
\qfigdef{p1_chem_s2_q09_B}{Score Advanced - Test 19 (04 Apr 2024)-img/p1_chem_s2_q09_B.png}{82.8}{0}
\qfigdef{p1_chem_s2_q09_C}{Score Advanced - Test 19 (04 Apr 2024)-img/p1_chem_s2_q09_C.png}{82.7}{0}
\qfigdef{p1_chem_s2_q09_D}{Score Advanced - Test 19 (04 Apr 2024)-img/p1_chem_s2_q09_D.png}{85.0}{0}
\qfigdef{p1_chem_s2_q10_1}{Score Advanced - Test 19 (04 Apr 2024)-img/p1_chem_s2_q10_1.png}{232.5}{0}
\qfigdef{p1_chem_s2_q10_A}{Score Advanced - Test 19 (04 Apr 2024)-img/p1_chem_s2_q10_A.png}{72.2}{0}
\qfigdef{p1_chem_s2_q10_B}{Score Advanced - Test 19 (04 Apr 2024)-img/p1_chem_s2_q10_B.png}{115.1}{0}
\qfigdef{p1_chem_s2_q10_C}{Score Advanced - Test 19 (04 Apr 2024)-img/p1_chem_s2_q10_C.png}{87.5}{0}
\qfigdef{p1_chem_s2_q10_D}{Score Advanced - Test 19 (04 Apr 2024)-img/p1_chem_s2_q10_D.png}{116.3}{0}
\qfigdef{p1_chem_s3_q07_1}{Score Advanced - Test 19 (04 Apr 2024)-img/p1_chem_s3_q07_1.png}{148.2}{0}
\qfigdef{p1_chem_s3_q08_1}{Score Advanced - Test 19 (04 Apr 2024)-img/p1_chem_s3_q08_1.png}{206.2}{0}
\raggedcolumns
\makeatletter\@beginparpenalty=10000\makeatother
\newcolumntype{L}[1]{>{\raggedright\arraybackslash}p{#1}}
\begin{document}
\begin{center}{\Large\bfseries Score Advanced --- Test 19 (04 Apr 2024)}\end{center}
\begin{multicols}{2}
\raggedcolumns
\papersec{1}
\partsec{1}{PHYSICS}
\sect{SECTION-I (i)}
\qn{1} A narrow monochromatic beam of light of intensity $I$ is incident on a glass plate labeled 1 as shown in figure (not drawn to scale). Another identical glass plate labeled 2 is kept close to the first one and parallel to it. Each glass plate reflects 25\% of the light incident on it and transmits the remaining. Find the ratio $\sqrt{\dfrac{I_{max}}{I_{min}}}$ of the interference pattern formed by two beams obtained after one reflection at each plate. ($I_{max}$ and $I_{min}$ represent the maximum value and minimum value of intensity respectively. Ignore the distance between plates 1 and 2 in the diagram. It is drawn for the sake of readability. Consider only the first reflection from the glass plates)\par
\qfig{p1_phy_s1_q01_1}
\begin{qoptions}
\opt{A}{3}
\opt{B}{5}
\opt{C}{6}
\opt{D}{7}
\end{qoptions}
\qn{2} One pan of a balance carries a container with water, and the other a stand with an iron sphere suspended from it (see the diagram). The pans are in balance as shown in figure a. Then the stand is turned around so the suspended sphere is completely submerged in the water. Obviously the balance disturbed, since the pan with the stand becomes lighter as shown in figure b. $w_1$ is the weight of the iron sphere $w_2$ is the weight of the water displaced by it when placed in water as shown in figure b. What additional approximate weight must be put on this pan to restore equilibrium? (Ignore the mass of the stand so that the torque effects are minimal)\par
\qfig{p1_phy_s1_q02_1}
\begin{qoptions}
\opt{A}{$w_1$}
\opt{B}{$w_2$}
\opt{C}{$2w_2$}
\opt{D}{$w_1+w_2$}
\end{qoptions}
\qn{3} A water surfer is moving at a speed of 15 m/s. When he is surfing in the direction of wave, he swings upwards every 0.8 s because of wave crests. While surfing in opposite direction to that of wave motion, he swings upwards every 0.6 s. Determine the wavelength of transverse component of the water wave.
\begin{qoptions}
\opt{A}{15m}
\opt{B}{10.3m}
\opt{C}{21.6m}
\opt{D}{25m}
\end{qoptions}
\qn{4} The following deuterium reactions and corresponding reaction energies are found to occur.\par
$^{14}$N(d, p)$^{15}$N,\quad Q = 8.53 MeV\par
$^{15}$N(d, $\alpha$)$^{13}$C,\quad Q = 7.58 MeV\par
$^{13}$C(d, $\alpha$)$^{11}$B,\quad Q = 5.16 MeV\par
The notation $^{14}$N(d, p)$^{15}$N represents the reaction $^{14}$N + d$\rightarrow$$^{15}$N + p\par
The masses of $^{4}_{2}$He = 4.0026 amu, $^{2}_{1}$H = 2.014 amu, $^{1}_{1}$H = 1.0078 amu, n=1.0087 amu (1amu = 931 MeV)\par
The Q values of the reaction $^{11}$B($\alpha$, n)$^{14}$N is
\begin{qoptions}
\opt{A}{0.5 eV}
\opt{B}{0.5 MeV}
\opt{C}{0.05 MeV}
\opt{D}{0.05 eV}
\end{qoptions}
\sect{SECTION-I (ii)}
\qn{5} The solar constant for the earth (s) is defined as energy received per unit area per second normally on the earth. The surface temperature of the sun is $T$ K. The sun subtends an angle $\theta$ at the earth. Then which of the following statements is/are correct:
\begin{qoptions}
\opt{A}{$s\propto T^4$}
\opt{B}{$s\propto T^2$}
\opt{C}{$s\propto\theta^2$}
\opt{D}{$s\propto\theta$}
\end{qoptions}
\qn{6} A straight conductor AB lies along the axis of hollow metal cylinder of length L. The cylinder is connected to earth through a conducting wire C. Some amount of charge will flow through C if one of the following changes is made from the given initial state. Mark the correct options for which charge flows through C:\par
\qfig{p1_phy_s2_q06_1}
\begin{qoptions}
\opt{A}{a current begins to flow through AB}
\opt{B}{a current that was already flowing through AB reverses its direction}
\opt{C}{AB is removed, and a beam of electron starts flowing in its place}
\opt{D}{AB is removed, and a magnet is dropped along AB with its north pole down}
\end{qoptions}
\qn{7} Three identical fixed metal plates are arranged in the air at equal distances $d$ from each other. The area of each of the plates is equal to $A$. The plate 1 has charge $+Q$, while plates 2 and 3 are not charged and are connected through a key K to the inductor L. The distance between the plates is small compared with their size. The ohmic resistance of the coil can be neglected.\par
\qfig{p1_phy_s2_q07_1}
\begin{qoptions}
\opt{A}{Just before the switch K is closed, the charge on right surface of plate 2 is $Q/2$}
\opt{B}{Just before the switch K is closed, the charge on right surface of plate 3 is $Q/2$}
\opt{C}{After the switch K is closed, the maximum current in the circuit is $\dfrac{Q}{2}\sqrt{\dfrac{d}{L\varepsilon_0A}}$}
\opt{D}{After the switch K is closed, the maximum current in the circuit is $Q\sqrt{\dfrac{d}{L\varepsilon_0A}}$}
\end{qoptions}
\qn{8} An ideal gas is taken from the state A ($P$, $V$) to the state B($P/2$, $2V$) along a straight line path as shown in figure. Select the correct statement from the following\par
\qfig{p1_phy_s2_q08_1}
\begin{qoptions}
\opt{A}{Work done by the gas in going from A to B exceeds the work done in going from A to B under isothermal conditions,}
\opt{B}{In the T-V diagram, part AB would becomes a parabola}
\opt{C}{In the P-T diagram, path AB would be part of hyperbola}
\opt{D}{In going from A to B, the temp. T of gas first increases to max value and then decreases.}
\end{qoptions}
\qn{9} Three spools A, B and C each having moment of inertia $I=\dfrac{MR^2}{4}$ are placed on rough ground and equal horizontal force $F$ is applied at positions as shown in the figures (a), (b) and (c). Then\par
\qfig{p1_phy_s2_q09_1}
\begin{qoptions}
\opt{A}{Friction force on spool A acts in forward direction}
\opt{B}{Friction force on spool A acts in backward direction}
\opt{C}{Friction force on spool B acts in forward direction}
\opt{D}{Friction force on spool B and C acts in backward direction}
\end{qoptions}
\qn{10} The figure shows a taut string, initially aligned with the x-axis, which carries a waveform travelling towards the right without any change its shape. The bold line and the dotted line indicate the displacement of the string (i.e., shape of the waveform) at time $=t$ second and at time $=(t+dt)$ second, respectively. The speed of particle P at $t=2$ sec is 1 cm/s. Choose the correct statement(s) from the following (assume that $dt$ is positive).\par
\qfig{p1_phy_s2_q10_1}
\begin{qoptions}
\opt{A}{\qfigi{p1_phy_s2_q10_A}}
\opt{B}{\qfigi{p1_phy_s2_q10_B}}
\opt{C}{Speed of waveform is 1 cm/s and it is moving along positive x-axis.}
\opt{D}{Speed of wavelength is 1 cm/s and it is moving along negative x-axis}
\end{qoptions}
\sect{SECTION-III}
\qn{1} A positive point charge $q=10e$ is fixed at the focus of an ellipse, with an electron revolving around it in the elliptical orbit. If the speed of electron at A is 2m/s then calculate the speed of electron at Point B. (eccentricity of ellipse = 0.2)\par
\qfig{p1_phy_s3_q01_1}
\qn{2} An infinite uniform current carrying wire is kept along z-axis, carrying current $I_0$ in the direction of the positive z-axis. OABCDEFG represents a circle (where all the points are equally spaced), whose centre is at point (4m, 0 m) and radius 4m as shown in the figure. If $\displaystyle\int_{DEF}\vec{B}\cdot d\vec{\ell}=\dfrac{\mu_0I_0}{k}$ in S.I unit, then the value of k is\par
\qfig{p1_phy_s3_q02_1}
\qn{3} The escape velocity for a planet is $v_e=\sqrt{6}$ km/s. A particle starts from rest at a large distance from the planet under the gravitational attraction of the planet and passes through a smooth tunnel dug along one of its diameters. The speed of the particle at the centre of planet will be (in km sec)
\qn{4} Two particles A and B of mass $m$ each are connected together by a rigid massless rod of length 20 cm. Initially rod is vertical and particle A is given velocity $V$ horizontally, while particle B is at rest. Consider the adjacent figure. Find the minimum value of $V$ (in m/s) for which particle B loses contact with ground immediately after being given velocity $V$. (Take $g=10$ m/s$^2$)\par
\qfig{p1_phy_s3_q04_1}
\qn{5} A cubical block of side length $l$ is allowed to fall by rotating clock wise from the rest position as shown in figure. The support point 'O' is a small frictionless roller. The angular velocity '$\omega$' of block as edge OC becomes horizontal just before striking the supporting surface is $\omega=\sqrt{\dfrac{6\sqrt2\left(2-\sqrt2\right)g}{kl}}$ where $k$ is a constant. Find the value of $k$.\par
\qfig{p1_phy_s3_q05_1}
\qn{6} The peak emission from a block body at a certain temperature occurs at a wavelength of 6200 \AA. On increasing its temperature, the total radiation emitted is increased 16 times. These radiations are allowed to fall on a metal surface. Photoelectrons emitted by the peak radiation at higher temperature can be bought to rest by applying a potential equivalent to the excited potential corresponding to the transition for the level $n=4$ to $n=2$ in the Bohr's hydrogen atom. The work function of the metal is given by $\dfrac{\alpha}{100}$eV where $\alpha$ is the numerical constant .Find the value of $\alpha$ and fill $\alpha/29$ .
\qn{7} A block was observed to move on a straight line uniformly. It displaces 10cm in 2 second. The straight line has calibrations of 1mm. and the stop watch used to measure time has least count of 0.1s. Using this data its speed is calculated Now the displacement of the block is calculated for next 10s using same stop watch. The absolute error in the calculated value of displacement will be $\dfrac{x}{2}$ cm then x is....
\qn{8} Many stars collapse under gravitational pull into a neutron star when most of the hydrogen is used up .The entire star is then a tightly packed ball of neutrons with the density of nucleus . Assume that a star of mass $2\times10^{32}$ kg, radius $5.5\times10^{8}$m collapses into a neutron star. It has a time period of rotation of $10^8$ sec about its axis initially .What is its rotation period (in sec) .after collapse .\par
(Take : $R_0=1.1$ fm, 1 amu $=1.6\times10^{-27}$ kg)
\partsec{2}{CHEMISTRY}
\sect{SECTION-I (i)}
\qn{1} Which of the following statement(s) is/are correct regarding complex ion $[\mathrm{Co(en)_2(NH_3)_2}]^{3+}$\par
(I) Central metal is sp$^3$d$^2$ hybridized\par
(II) Complex ion can show geometrical as well as optical isomerism.\par
(III) It has two N$-$Co$-$N bond angles.
\begin{qoptions}
\opt{A}{I \& II}
\opt{B}{II \& III}
\opt{C}{Only II}
\opt{D}{Only III}
\end{qoptions}
\qn{2} Identify the CORRECT statements.\par
(I) NH$_4$Cl is added along with NH$_4$OH so that only III$^{rd}$ group cations can be precipitated as their hydroxides.\par
(II) In III$^{rd}$ group in place of NH$_4$Cl, (NH$_4$)$_2$SO$_4$ cannot be used as V$^{th}$ group radical(s) will also be precipitated as MSO$_4$ along with Al$^{3+}$, Fe$^{3+}$ \& Cr$^{3+}$.\par
(III) Aqueous solution of ammonium sulphate produces white ppt. with BaCl$_2$ solution which is insoluble in dil.HCl.
\begin{qoptions}
\opt{A}{I and II only}
\opt{B}{II and III only}
\opt{C}{I and III only}
\opt{D}{I, II and III}
\end{qoptions}
\qn{3} The major product of the following reaction is\par
\qfig{p1_chem_s1_q03_1}
\begin{qoptions}
\opt{A}{\qfigi{p1_chem_s1_q03_A}}
\opt{B}{\qfigi{p1_chem_s1_q03_B}}
\opt{C}{\qfigi{p1_chem_s1_q03_C}}
\opt{D}{\qfigi{p1_chem_s1_q03_D}}
\end{qoptions}
\qn{4}\par
\qfig{p1_chem_s1_q04_1}
The product [M] is :
\begin{qoptions}
\opt{A}{\qfigi{p1_chem_s1_q04_A}}
\opt{B}{\qfigi{p1_chem_s1_q04_B}}
\opt{C}{\qfigi{p1_chem_s1_q04_C}}
\opt{D}{\qfigi{p1_chem_s1_q04_D}}
\end{qoptions}
\sect{SECTION-I (ii)}
\qn{5} In wave mechanical model, wave function of orbital is represented by $\psi_{n,\ell,m_\ell}$. Based on the following given wave function of orbitals\par
$\psi_{1,0,0}=\dfrac{2}{a_0^{3/2}}e^{-r/na_0}$\par
$\psi_{3,0,0}=\dfrac{1}{9\sqrt3}\left(\dfrac{1}{a_o}\right)^{3/2}\left(6-6\sigma+\sigma^2\right)e^{-\sigma/2}$ Where $\sigma=\dfrac{2r}{na_0}$\par
$\psi_{2,0,0}=\dfrac{1}{2\sqrt2}\left(\dfrac{1}{a_0}\right)^{3/2}\left(2-\dfrac{2r}{na_0}\right)e^{-r/na_0}$\par
$\psi_{3,1,0}=\dfrac{1}{9\sqrt6}\left(\dfrac{1}{a_0}\right)^{3/2}(4-\sigma)\sigma e^{-\sigma/2}$, where $\sigma=\dfrac{2r}{na_0}$\par
Select the correct statement(s)
\begin{qoptions}
\opt{A}{For 3p orbital, the graph will intersect r-axis at two distinct points.}
\opt{B}{For 1s orbital of H-atom, the radial probability distribution of finding electron will be maximum at $r=a_0$.}
\opt{C}{For 3s orbital, number of angular nodes is zero.}
\opt{D}{For 2s orbital, radial distance for radial node is $2a_0$}
\end{qoptions}
\qn{6} Potassium crystallizes in a body centered cubic lattice with edge length $a=5.0$\AA. Then choose the correct statements in given option -
\begin{qoptions}
\opt{A}{The number of nearest neighbours to any K atom is 8}
\opt{B}{The number of next nearest neighbours to any K atom is 6}
\opt{C}{Distance between nearest neighbours is 4.33 \AA}
\opt{D}{Distance between next nearest neighbours is 5.0\AA}
\end{qoptions}
\qfullbegin
\qn{7} Select correct for the following given two types of double chain silicates anion\par
\qfig{p1_chem_s2_q07_1}
\qfullsplit
\begin{qoptions}
\opt{A}{each Si in both is bonded with four 'O' atoms}
\opt{B}{oxidation number of each Si $=+4$ in both}
\opt{C}{both have tetrahedron with three shared 'O' corners}
\opt{D}{in both all the tetrahedron have three shared 'O' corners}
\end{qoptions}
\qfullend
\qn{8} A compound A was isolated as a major product formed by the reaction between a s-block metal M and excess O$_2$. A reacts with water to give only MOH. while M reacts in a controlled manner with water giving MOH and another product, B. Select CORRECT options for metal "M".
\begin{qoptions}
\opt{A}{Chloride of "M" gives red colour on flame.}
\opt{B}{A compound "C" is formed on reaction with metal M and nitrogen, produce a basic gas on hydrolysis.}
\opt{C}{Carbonate and bicarbonates of metal "M" are crystalline solid.}
\opt{D}{Chloride of metal M is deliquescent and crystalize as hydrated MCl.2H$_2$O}
\end{qoptions}
\qn{9} Compound I reacts with metallic sodium in dry Et$_2$O to form compound J, which on heating forms K (C$_4$H$_6$). 1 mole of hydrogen bromide reacts with K to give L (C$_4$H$_7$Br) and M (C$_4$H$_7$Br) in ratio 80 : 20 respectively.\par
\qfig{p1_chem_s2_q09_1}
\begin{qoptions}
\opt{A}{\qfigi{p1_chem_s2_q09_A}}
\opt{B}{\qfigi{p1_chem_s2_q09_B}}
\opt{C}{\qfigi{p1_chem_s2_q09_C}}
\opt{D}{\qfigi{p1_chem_s2_q09_D}}
\end{qoptions}
\qn{10}\par
\qfig{p1_chem_s2_q10_1}
\begin{qoptions}
\opt{A}{\qfigi{p1_chem_s2_q10_A}}
\opt{B}{\qfigi{p1_chem_s2_q10_B}}
\opt{C}{\qfigi{p1_chem_s2_q10_C}}
\opt{D}{\qfigi{p1_chem_s2_q10_D}}
\end{qoptions}
\sect{SECTION-III}
\qn{1} Ozone acts as powerful oxidizing agent. When ozone reacts with PbS, I$^{-}$ \& NO, libirates O$_2$.\par
$x$ moles of O$_3$ are required to oxidize each mole of PbS\par
$y$ moles of O$_3$ are required to oxidize each mole of I$^{-}$ in basic medium\par
$z$ moles of O$_3$ are required to oxidize each mole of NO\par
Find $(x+y+z)$.
\qn{2} For the reaction: Br$_2$($\ell$) + Cl$_2$(g) $\rightleftharpoons$ 2BrCl(g);$K_p=1$ atm at 300K. In a closed container of volume 164L. initially 10 moles of Cl$_2$(g) are present. If '$m$' gm is the minimum mass of Br$_2$($\ell$) which must be introduced into this container so that above equilibrium is maintained at total pressure of 2.25 atm, then the value of $\left(\dfrac{m}{100}\right)$ is.\par
[Given: vapour pressure of Br$_2$($\ell$) at 27$^\circ$C $=0.25$ atm, volume of liquid Br$_2$, is negligible but sufficient to achieve equilibrium. Temperature of 300K is maintained throughout;\par
$R=0.082$ L-atm / K/ mol; Atomic mass of Br $=80$]\par
Fill your answer as sum of digits (excluding decimal places) till you get the single digit answer.
\qn{3} The saturated solution of a certain salt show elevation in boiling point equal to $3.12\times10^{-3}$ K in aqueous solution. If the solubility of salt is $2\times10^{-3}$ mole/kg at boiling point, then how many moles of ions are given by 1 mole of solid salt on complete dissociation. ($K_b$(H$_2$O) $=0.52$ K kg/mole)
\qn{4} 0.01 moles of sulphuryl chloride SO$_2$Cl$_2$(g) is taken in a sealed tube \& heated to 400 K.\par
It decomposes by following 1st order kinetics according to reaction:\par
SO$_2$Cl$_2$(g) $\rightarrow$ SO$_2$(g) + Cl$_2$(g)\par
The tube is broken after 4 hours \& gas passed through a 25 ml of an acidified 1 N iodine solution where all SO$_2$ is oxidized to SO$_4^{2-}$ (consider other components are not involved in reaction with iodine). The resulting solution required 10 ml of 1 M Hypo solution. What is half life in hours.
\qn{5} Resistance of a conductivity cell filled with 0.01 M HA (aq.) solution is 200$\Omega$ at 300K. (HA is weak monoprotic acid). If cell constant is 4 cm$^{-1}$ \& $\Lambda_m^0$(HA) is 2$\Omega^{-1}$m$^2$mol$^{-1}$. Find the pH of the given acid solution.
\qn{6} Number of negatively charged colloids among the following is -\par
(i) As$_2$S$_3$ sol\quad (ii) Ferric hydroxide sol\par
(iii) Gold sol\quad (iv) Clay\par
(v) Acidic dye\quad (vi) Hemoglobin\par
(vii) Basic dye\quad (viii) Starch
\qn{7} Observe the following reaction and find out that how many number of reactant stereoisomers can be reduced to optically inactive meso product\par
\qfig{p1_chem_s3_q07_1}
\qn{8}\par
\qfig{p1_chem_s3_q08_1}
Find the number of lone pair of electron in product [T].
\partsec{3}{MATHEMATICS}
\sect{SECTION-I (i)}
\qn{1} Let A, B, C and D be four points in space whose coordinates $(x, y, z)$ satisfy $x^4+y^4+z^4+1=4xyz$, then the volume of tetrahedron ABCD is
\begin{qoptions}
\opt{A}{16}
\opt{B}{$\dfrac{16}{3}$}
\opt{C}{$\dfrac{8}{3}$}
\opt{D}{$\dfrac{4}{3}$}
\end{qoptions}
\qn{2} Given that $P(A)=2/7$, $P(B)=2/7$, $P\left(A\cap\overline{B}\cap C\right)=1/14$, $P\left(A\cap\overline{B}\cap\overline{C}\right)=1/14$. If $P\left(\overline{A}|B\right)=\dfrac{x}{y}$ where $x$, $y$ are coprime positive integers then which of the following statements is true?
\begin{qoptions}
\opt{A}{$x+y=7$}
\opt{B}{$x+y=9$}
\opt{C}{$x+y=13$}
\opt{D}{$x+y=15$}
\end{qoptions}
\qn{3} $\displaystyle\lim_{a\to\infty}a^2\int_0^{1/a^2}\dfrac{\cos^{-1}ax}{\sec^{-1}(ax+2)}dx$ is equal to
\begin{qoptions}
\opt{A}{3}
\opt{B}{$\dfrac32$}
\opt{C}{$\dfrac23$}
\opt{D}{2}
\end{qoptions}
\qn{4} If $r$, $s$, $t$ are the prime numbers and $p$, $q$ are the positive integers such that the LCM of $p$ \& $q$ is $r^2t^4s^2$ then the number of ordered pair $(p, q)$ is:
\begin{qoptions}
\opt{A}{252}
\opt{B}{254}
\opt{C}{225}
\opt{D}{224}
\end{qoptions}
\sect{SECTION-I (ii)}
\qn{5} If $x^2-(a-3)x+a=0$ has at least one positive root then $a$ may belong to
\begin{qoptions}
\opt{A}{$(-\infty,0)$}
\opt{B}{$(-\infty,0)\cup[7,\infty)$}
\opt{C}{$[9,\infty)$}
\opt{D}{$(-\infty,0)\cup[10,\infty]$}
\end{qoptions}
\qn{6} 3 circles of radii 1, 2 and 3 and centres at A, B and C respectively, touch each other. Another circle whose centre is P touches all these 3 circles externally and has radius $r$. Also $\angle PAB=\theta$ and $\angle PAC=\alpha$, then
\begin{qoptions}
\opt{A}{$\cos\theta=\dfrac{3-r}{3(1+r)}$}
\opt{B}{$\cos\alpha=\dfrac{2-r}{2(1+r)}$}
\opt{C}{$r=\dfrac{6}{23}$}
\opt{D}{$r=\dfrac{6}{\sqrt{23}}$}
\end{qoptions}
\qn{7} Let C be a curve such that normal at any point P on it meets x - axis and y - axis at A and B respectively. If BP : PA = 1 : 2 (internally) and curve passes through $(0, 4)$, then -
\begin{qoptions}
\opt{A}{The curve 'C' passes through $\left(\sqrt{10},-6\right)$}
\opt{B}{Equation of tangent at $\left(4,4\sqrt3\right)$ is $\sqrt3x-2y+4\sqrt3=0$}
\opt{C}{Equation of tangent at $\left(4,4\sqrt3\right)$ is $2x-\sqrt3y+4=0$}
\opt{D}{If point P is $(8,12)$ then point A and point B are $(24, 0)$ and $(0, 18)$ respectively.}
\end{qoptions}
\qn{8} For all real number $x$, $2^x+3^x-4^x+6^x-9^x$ is always less than -
\begin{qoptions}
\opt{A}{0}
\opt{B}{1}
\opt{C}{2}
\opt{D}{3}
\end{qoptions}
\qn{9} Let A denote the area bounded by $y=\ell n^4x$ and lines $y=0$ and $x=e^4$, then -
\begin{qoptions}
\opt{A}{$A<256e^4-175e^3-81$}
\opt{B}{$A<256e^4-256$}
\opt{C}{$A>256\left(e^4-1\right)$}
\opt{D}{$A>\dfrac{337}{2}e^3(e-1)$}
\end{qoptions}
\qn{10} A is a matrix of order $3\times2$ and B is a matrix of order $2\times3$. If $AB=C$ and $BA=D$ then
\begin{qoptions}
\opt{A}{$\det(C)=0$}
\opt{B}{$\det(C)\cdot\det(D)\neq0$}
\opt{C}{$\det(C)\cdot\det(D)=0$}
\opt{D}{$\det(C)\neq0$}
\end{qoptions}
\sect{SECTION-III}
\qn{1} Let A and B be $3\times3$ square matrices such that $AB=9I$ where $A=\begin{pmatrix}0&1&-1\\4&-3&4\\3&-3&\lambda\end{pmatrix}$. If $B_{33}=9A_{23}$, then the value of $\mathrm{Tr}\left(\dfrac{A^2+B^2}{82}\right)+\lambda$, is\par
(where $x_{ij}$ represents elements of i$^{th}$ row and j$^{th}$ column of square matrix $x$ and Tr($x$) represents trace of square matrix $x$).
\qn{2} A \& B are two separate reservoirs of water. Capacity of reservoir A is double the capacity of reservoir B. Both the reservoirs are filled completely with water, their inlet are closed and then the water is released simultaneously from both the reservoirs. The rate of flow of water out of each reservoir at any instant of time is proportional to the quantity of water in the reservoir at that time. One hour after the water is released, the quantity of water in reservoir A is 1.5 times the quantity of water in reservoir B.\par
After $\log_a2$ hours both the reservoirs have the same quantity of water then $\dfrac4a$ is equal to
\qn{3} Tangents are drawn from any point on the hyperbola $\dfrac{x^2}{9}-\dfrac{y^2}{4}=1$ to the circle $x^2+y^2=9$. If the locus of the mid-point of the chord of contact is $\alpha(x^4+y^4+2x^2y^2)=\beta x^2-\gamma y^2$ then the value of $(\gamma-2\beta-\alpha)$ is equal to
\qn{4} If the value of $\displaystyle\lim_{x\to\infty}x^4\left(\left(\dfrac{x^2+1}{x^2-1}\right)^{x^2}-e^2\right)$ is equal to $\dfrac{ae^2}{b}$ (where $a$ and $b$ are co - prime) then the value of $\left|a^2-b^2\right|$ is
\qn{5} If $\dfrac{1}{\cot9^o-3\tan9^o}+\dfrac{3}{\cot27^o-3\tan27^o}+\dfrac{9}{\cot81^o-3\tan81^o}+\dfrac{27}{\cot243^o-3\tan243^o}=A\tan B^o$ where A \& B are positive integers and $B<50$, find $|A-B|$.
\qn{6} Consider a twice differentiable function $y=f(x)$ satisfying $f(x)+f''(x)=2f'(x)$ where $f(0)=0$, $f(1)=e$ then find $\displaystyle\int_0^1f(x)dx$
\qn{7} Number of solutions of $x\in[0,\pi]$ satisfying the equation\par
$\left(\log_{\sqrt3}\tan x\right)\left(\sqrt{\dfrac{\log_{\sqrt3}\left(\sqrt3\right)^3}{\log_{\sqrt3}\sqrt3}+\dfrac{\log_{\sqrt3}3}{\log_{\sqrt3}\tan x}}\right)=-1$ is
\qn{8} The minimum distance of the point $(1,1,1)$ from $x+y+z=1$ measured perpendicular to the line $\dfrac{x-x_1}{1}=\dfrac{y-y_1}{2}=\dfrac{z-z_1}{3}$ is $k$ then $\dfrac{9k}{\sqrt{21}}=$ \underline{\hspace{1.5cm}}
\end{multicols}
\end{document}
